\documentclass[12pt]{article}
\usepackage[utf8]{inputenc}
\usepackage[T1]{fontenc}
\usepackage[margin=1in]{geometry}
\usepackage{setspace}
\usepackage{xcolor}
\usepackage[colorlinks=true, urlcolor=blue!60!black, linkcolor=blue!60!black]{hyperref}
\setlength{\parindent}{2em}
\setlength{\parskip}{1.25em}
\renewcommand{\baselinestretch}{1.0}
\begin{document}
\begin{center}
{\LARGE \textit{Pike v. Bruce Church, Inc.}, 397 U.S. 137 (1970)\par}
\vspace{0.8em}
{\large Case Brief --- Four-Page Edition\par}
\end{center}
\vspace{1em}

\section{Memory Jogger}

Arizona ordered a grower to pack its premium Parker cantaloupes in Arizona rather than in its California sheds 31 miles away, so the melons would carry Arizona's name; a unanimous Court struck the order under the \textit{Pike} balancing test, which invalidates evenhanded laws whose burden on interstate commerce is clearly excessive in relation to their putative local benefits (\href{https://supreme.justia.com/cases/federal/us/397/137/}{397 U.S. 137}).

\section{Detailed Case Facts}

Bruce Church, Inc., farmed and packed produce in Arizona and California for interstate shipment. Under a 1964 lease with the Colorado River Indian Tribes and federal officials, it invested more than \$3,000,000 to develop about 6,400 acres of arid land at Parker, Arizona, and from 1966 grew cantaloupes regarded as higher in quality than others grown in Arizona. Cantaloupes are highly perishable and must be packed immediately in specialized sheds. Lacking sheds at Parker, the company trucked its melons in bulk 31 miles to its efficient packing sheds in Blythe, California, where they were packed under California standards virtually identical to Arizona's, in crates bearing the California packer's name.

In 1968, Loren J. Pike, the official enforcing Arizona's Fruit and Vegetable Standardization Act, invoked a provision requiring Arizona-grown cantaloupes to ``be packed in regular compact arrangement in closed standard containers approved by the supervisor'' and ordered the company not to ship uncrated melons out of state. With no Arizona shed nearby, the company faced the loss of a 1968 crop worth \$700,000.

\section{Procedural History}

Bruce Church sued to enjoin the order as an unconstitutional burden on interstate commerce. A three-judge district court granted preliminary relief, declining to abstain given the emergency. The parties stipulated that compliance would require building packing facilities near Parker over many months at a cost of about \$200,000. The court issued a permanent injunction, and Pike appealed directly; the Supreme Court affirmed on March 2, 1970.

\section{Judicial Votes}

Affirmed, \textbf{8--0}, with one seat vacant.
\begin{itemize}
\item \textbf{Majority:} Justice Stewart, joined by Chief Justice Burger and Justices Black, Douglas, Harlan, Brennan, White, and Marshall.
\item \textbf{Concurrences and dissents:} None.
\end{itemize}

\section{Holding}

The order burdens interstate commerce because the melons were destined for immediate out-of-state shipment and the order required an out-of-state operation to be performed in Arizona. Where a statute regulates evenhandedly to effectuate a legitimate local interest with only incidental interstate effects, it stands unless the burden is ``clearly excessive in relation to the putative local benefits.'' Arizona's minimal interest in identifying the origin of superior melons to enhance other growers' reputations cannot justify forcing the company to build an unneeded \$200,000 plant.

\section{Analysis of Opinions}

\textbf{Majority (Stewart): commerce nexus.} Stewart distinguished the Mississippi cotton-tax cases, where the cotton had come to rest with ``no ascertainable destination without the state.'' Arizona instead demanded that an out-of-state operation be relocated in-state, a mandate invalidated in \textit{Foster-Fountain Packing Co. v. Haydel}, \textit{Johnson v. Haydel}, and \textit{Toomer v. Witsell}.

\textbf{The balancing framework.} Synthesizing \textit{Huron Portland Cement} and \textit{Southern Pacific}, the Court asked whether the law regulates evenhandedly with only incidental effects, whether the burden is clearly excessive, and whether the local interest ``could be promoted as well with a lesser impact on interstate activities.''

\textbf{Application.} The Act generally served legitimate ends, such as preventing deceptive packaging and contaminated produce, but this order involved no health, safety, or deception: the superior melons were properly packed and not labeled as Arizona produce. Arizona sought only ``reflected good will,'' an interest ``tenuous'' and ``minimal, at best.'' The burden's nature mattered more than its extent: forcing operations in-state that could be done more efficiently elsewhere is ``virtually per se illegal.'' If a state cannot require local processing to secure jobs, it cannot impose a ``straitjacket'' on a company to boost others' reputations.

\textbf{Later developments.} In \textit{National Pork Producers Council v. Ross} (2023), the Court rejected an ``almost per se'' extraterritoriality rule and fractured over \textit{Pike}: Justice Gorsuch's majority recast it as a tool to ``smoke out'' protectionism; three Justices found economic costs and moral interests incommensurable; and six Justices preserved \textit{Pike} balancing.

\section{Examples: Future Applications}

\begin{enumerate}
\item \textbf{Same-Side (Invalid): In-State Data Servers.} A state requires spinach to be logged on in-state servers, forcing a \$450,000 redundant data center despite instant remote access. Mandating local operations still fails.
\item \textbf{Same-Side (Invalid): Shrimp Processing.} A state requires its shrimp to be peeled in-state before export, as in \textit{Foster-Fountain}, a burden \textit{Pike} called virtually per se illegal.
\item \textbf{Opposite-Side (Valid): Packaging Ban.} An evenhanded ban on polystyrene clamshells imposes \$1.2 million in retooling but no local-processing mandate (\textit{Clover Leaf Creamery}, \textit{Exxon}, \textit{Ross}).
\item \textbf{Opposite-Side (Valid): Smoke Ordinance.} A port city's smoke limits applied to docked ships regulate evenhandedly for health, as in \textit{Huron Portland Cement}.
\item \textbf{Fence-Sitter: Hormone-Free Milk.} A ban on milk from hormone-treated cows forces out-of-state cooperatives to build segregated supply chains; given \textit{Ross}'s divisions, the balance is uncertain.
\end{enumerate}

\section{Critique}

\textbf{Hidden protectionism versus true balancing.} Although famed for balancing, \textit{Pike} relied heavily on the ``virtually per se illegal'' rule against forced local processing, which supports \textit{Ross}'s rereading of the case as aimed at latent protectionism of local packers.

\textbf{Incommensurability and unexplored alternatives.} \textit{Pike} offers no metric for weighing compliance dollars against noneconomic values, as \textit{Bendix} later noted, and it omitted an obvious less restrictive alternative: requiring origin stamps on California-packed crates.

\section{Key Quotations}

\begin{quote}
``Although the criteria for determining the validity of state statutes affecting interstate commerce have been variously stated, the general rule that emerges can be phrased as follows: where the statute regulates evenhandedly to effectuate a legitimate local public interest, and its effects on interstate commerce are only incidental, it will be upheld unless the burden imposed on such commerce is clearly excessive in relation to the putative local benefits. \textit{Huron Cement Co. v. Detroit}, 362 U. S. 440, 443. If a legitimate local purpose is found, then the question becomes one of degree. And the extent of the burden that will be tolerated will, of course, depend on the nature of the local interest involved, and on whether it could be promoted as well with a lesser impact on interstate activities.''
\end{quote}
\noindent --- Justice Stewart, for the Court, 397 U.S. at 142.

\begin{quote}
``The nature of that burden is, constitutionally, more significant than its extent. For the Court has viewed with particular suspicion state statutes requiring business operations to be performed in the home State that could more efficiently be performed elsewhere. Even where the State is pursuing a clearly legitimate local interest, this particular burden on commerce has been declared to be virtually per se illegal.''
\end{quote}
\noindent --- Justice Stewart, for the Court, 397 U.S. at 145.

\begin{quote}
``While the order issued under the Arizona statute does not impose such rigidity on an entire industry, it does impose just such a straitjacket on the appellee company with respect to the allocation of its interstate resources. Such an incidental consequence of a regulatory scheme could perhaps be tolerated if a more compelling state interest were involved. But here, the State's interest is minimal, at best -- certainly less substantial than a State's interest in securing employment for its people.''
\end{quote}
\noindent --- Justice Stewart, for the Court, 397 U.S. at 146.

\end{document}
